\section{问题重述}

本模板以一类典型的配送路径优化问题作为示意场景，用于展示正式论文在章节层次、公式编号、图表位置和参考文献引用方面的排版效果。后续在替换为真实赛题时，可直接保留这一结构，只修改问题背景、数据来源和研究目标等内容。\cite{dantzig_ramser_1959,solomon_1987_vrptw}

\section{问题分析}

从建模角度看，此类问题通常需要同时处理运输成本、服务时效和调度稳定性等多个目标，并在容量、时窗、路网状态等约束下构建可求解的数学模型。因此，正文结构宜先完成问题拆解，再进入变量定义、模型建立、求解方法和结果解释。

\subsection{任务结构拆解}

在论文写作中，可以先将总问题拆解为若干彼此相关但边界清晰的子问题，例如订单分配、路径排序和结果评估。这样既有利于控制每一节的叙述长度，也便于在后续补充图表和实验结果时维持整体逻辑的一致性。

\subsubsection{目标层次}

如果问题同时包含经济性与服务性目标，可先给出各目标的业务含义，再说明采用加权求和、层次优化或约束转化的方式统一建模。对比赛论文而言，这种写法通常比直接堆叠公式更容易建立评审对模型设计的整体理解。

\subsubsection{约束层次}

约束条件建议按资源约束、时间约束和业务规则约束三个层面分组描述。这样的写法有助于后续在模型建立章节中逐条落到变量和公式上，也更方便在答辩或复审时定位假设来源。
